\section{拉格朗日子流形的经典形变}
\label{sec:classical-lagrangian-deformations}
\stfhSubIndexRangeStart{I186}{la^ge^lang^ri^zi^liu^xing}{拉格朗日子流形（Lagrangian submanifold）}{jing^dian^xing^bian}{经典形变（classical deformations）}

达布--温斯坦定理的一个直接推论，完整描述了给定拉格朗日子流形附近的拉格朗日子流形之局部参数化。下面把这一陈述说得更精确。

先引入一个一般定义。

\begin{definition}\label{def:hamiltonian-isotopic-lagrangians}
\stfhSubIndex{I154}{ha^mi^dun}{哈密顿（Hamiltonian）}{ha^mi^dun^tong^hen}{哈密顿同痕（Hamiltonian isotopy）}
若存在含时哈密顿量 $H:[0,1]\times M\to\mathbb R$ 使
\[
 L_1=\phi_H^1(L_0),
\]
则称两个拉格朗日子流形 $L_0,L_1$ \emph{哈密顿同痕}。称同痕
$\{L_t=\phi_H^t(L_0)\}$ 为 $L_0$ 的一个\emph{哈密顿同痕}。
\end{definition}

\subsection{零截面的情形}
\label{subsec:deformations-zero-section}

回顾：一个 $1$-形式（也即余切丛 $T^*M\to M$ 的一个截面）的图像是拉格朗日的，当且仅当这个 $1$-形式闭。此外，任何在 $C^1$ 意义下充分靠近零截面
$o_L\subset T^*L$ 的拉格朗日子流形都是某个 $1$-形式的图像。特别地，如果把 $T^*L$ 的拉格朗日子流形集合
$\operatorname{Lag}(T^*L)$ 看作无限维（弗雷歇）流形，那么零截面的邻域以闭 $1$-形式集合 $Z^1(L)$ 为模型。进一步，任何在 $C^1$ 意义下靠近零截面的\emph{正合}拉格朗日子流形都与零截面哈密顿同痕。这是因为任一这样的拉格朗日子流形都形如
$\operatorname{Graph}(df)$，其中函数 $f$ 在 $C^2$ 意义下靠近常值函数；此时同痕
\[
 t\in[0,1]\longmapsto\operatorname{Graph}(t\,df)
\]
就是所需的哈密顿同痕。

\begin{exercise}\label{ex:zero-section-graph-hamiltonian}
证明上述零截面的同痕由（自治）函数
$H=f\circ\pi:T^*L\to\mathbb R$ 诱导。
\end{exercise}

\begin{exercise}\label{ex:contractible-neighborhood-symp-ham}
证明 $\operatorname{Symp}(M,\omega)$（相应地，$\operatorname{Ham}(M,\omega)$）中恒等映射有一个可缩的 $C^\infty$ 邻域。（提示：考虑辛微分同胚
$\phi:M\to M$ 与其图像
$\operatorname{Graph}\phi\subset M\times M$ 之间的对应；后者关于 $M\times M$ 上的形式
$(-\omega)\oplus\omega$ 是拉格朗日子流形。）
\end{exercise}

更一般地，在余切丛上，拉格朗日子流形模哈密顿同痕有如下局部描述。

\begin{proposition}\label{prop:local-lagrangians-mod-hamiltonian}
在 $C^1$ 意义下靠近零截面的拉格朗日子流形集合，模去 $C^1$-小的哈密顿同痕后，微分同胚于 $H^1(L,\mathbb R)$ 中零点的一个开邻域。
\end{proposition}

\begin{proof}
设 $\mathcal N$ 是 $\operatorname{Lag}(T^*L)$ 中零截面的一个 $C^1$ 邻域。如上所述，对某个 $\delta>0$，可写成
\begin{equation}
 \mathcal N
 =\{\operatorname{Graph}\alpha\mid\alpha\in Z^1(L),\ |\alpha|_{C^1}<\delta\}.
 \label{eq:lagrangian-zero-neighborhood}
\end{equation}
以 $L_\alpha$ 表示与 $\alpha$ 对应的拉格朗日子流形。考虑映射
\[
 L_\alpha\in\mathcal N\longmapsto[\alpha]\in H^1(L,\mathbb R).
\]
显然，该映射满射到 $H^1(L,\mathbb R)$ 中零点的一个开邻域。

假设 $L_\alpha$ 与 $L_{\alpha'}$哈密顿同痕，即对某个函数
$H:[0,1]\times M\to\mathbb R$ 有
$L_{\alpha'}=\phi_H^1(L_\alpha)$，并且对所有 $t\in[0,1]$，
$\phi_H^t(L_\alpha)\subset\mathcal N$。于是可写
\begin{equation}
 \phi_H^t(L_\alpha)=L_{\alpha_t}
 \label{eq:hamiltonian-graph-family}
\end{equation}
其中 $\alpha_t$ 是一族闭 $1$-形式。为计算上同调类 $[\alpha_t]$，对环路
$c:S^1\to L$ 计算 $\int_c\alpha_t$。把 $1$-形式 $\alpha,\alpha_t$ 等同于映射
$\widetilde\alpha,\widetilde\alpha_t:L\to T^*L$，则
\[
 \widetilde\alpha_t=\phi_H^t\circ\widetilde\alpha.
\]
因此，回顾 $\beta=(\widetilde\beta)^*\theta$，有
\[
 \int_c\alpha_t
 =\int_c(\widetilde\alpha_t)^*\theta
 =\int_c(\widetilde\alpha)^*(\phi_H^t)^*\theta
 =\int_{\widetilde\alpha\circ c}(\phi_H^t)^*\theta.
\]
计算
\[
\begin{aligned}
 \int_c\alpha'-\int_c\alpha
 &=\int_0^1\left(\int_c\frac d{dt}\alpha_t\right)dt\\
 &=\int_0^1\left(\int_{\widetilde\alpha\circ c}
     \frac d{dt}(\phi_H^t)^*\theta\right)dt.
\end{aligned}
\]
在上式中把\eqref{eq:liouville-form-under-hamiltonian-flow}用于
$\alpha=-\theta$，得到
\[
 \int_{\widetilde\alpha\circ c}\frac d{dt}(\phi_H^t)^*\theta=0.
\]
故对任意环路 $c$，$\int_c\alpha'-\int_c\alpha=0$。这证明
$[\alpha]=[\alpha']\in H^1(L;\mathbb R)$，因而上述赋值诱导良定义的满射
\[
 \mathcal N/\operatorname{Ham}\longrightarrow H^1(L,\mathbb R),
\]
其像是 $H^1(L,\mathbb R)$ 中零点的一个开邻域。还须证明单射性。设
$[\alpha]=[\alpha']$，亦即 $\alpha'-\alpha=df$。与
\cref{ex:zero-section-graph-hamiltonian}类似，单射性由下面的
\cref{ex:closed-form-graph-hamiltonian}得到。这就完成了证明。
\end{proof}

\begin{exercise}\label{ex:closed-form-graph-hamiltonian}
证明映射
\[
 t\longmapsto\alpha_t:=\alpha+t\,df
\]
以及与它相伴的拉格朗日子流形 $L_{\alpha_t}$，是由函数
$H=f\circ\pi$ 所相伴的哈密顿同痕诱导的。
\end{exercise}

\subsection{一般情形}
\label{subsec:general-lagrangian-deformation}

现在考虑一般的 $(M,\omega)$。设 $L\subset M$ 是拉格朗日嵌入。达布--温斯坦定理给出微分同胚
\stfhIndex{I086}{yu^qie^cong^zuo^biao^ka}{余切丛坐标卡（cotangent bundle chart）}
\[
 \Phi:U\subset M\longrightarrow V\subset T^*L,
\]
它延拓映射 $L\to o_L$，并且 $\Phi^*\omega_0=\omega$。称 $\Phi$ 为 $L$ 的\emph{余切丛坐标卡}。

注意，所有在 $C^1$ 意义下充分靠近给定拉格朗日子流形 $L\subset M$ 的拉格朗日子流形（事实上，这一目的下 $C^0$-接近就够了）都包含在一个给定达布坐标卡中；该坐标卡 $\Phi$ 把 $L$ 映到
$o_L\subset T^*L$。由上面对余切丛的讨论可知，任一这样的、在 $C^1$ 意义下靠近 $L\subset M$ 的拉格朗日子流形，都可由 $L$ 上的闭 $1$-形式参数化；这些形式位于 $Z^1(L)$ 中零点的一个开 $C^1$ 邻域。需要强调，这种参数化依赖余切丛坐标卡，但它的无穷小版本并不依赖。

先研究两个余切丛坐标卡 $\Phi_1,\Phi_2$ 的转移映射；两者均延拓典范等同
$Z:L\to o_L\subset T^*L$。令
\[
 \Phi_{12}=\Phi_2\circ(\Phi_1)^{-1}:V\longrightarrow V'.
\]
它满足 $\Phi_{12}^*\omega_0=\omega_0$ 以及
$\Phi_{12}|_{o_L}=\operatorname{id}|_{o_L}$。回顾：切空间
$T_L\operatorname{Sub}(M)$ 典范同构于 $\Gamma(\nu(L))$，其中
$\nu(L)=TM/TL$ 是 $L$ 的法丛，而 $\operatorname{Sub}(M)$ 是 $M$ 的 $n$ 维子流形集合。

\begin{lemma}\label{lem:normal-cotangent-identification}
有典范同构
\[
 \Gamma(\nu(o_L))\cong\Gamma(T^*L),
\]
它由
$[X]\mapsto(X\mathbin{\lrcorner}\omega_0)|_{To_L}$ 给出，其中 $X$ 是沿
$o_L$ 的向量场。在这个等同下，有自然等同
\begin{equation}
 T_{o_L}\operatorname{Lag}(T^*L)\cong Z^1(L),
 \label{eq:tangent-lag-closed-forms}
\end{equation}
并且 $d\Phi_{12}$ 在 $\Gamma(T^*L)$ 上诱导恒等映射。
\end{lemma}

\begin{proof}
为证明第一项，回顾典范分裂
$T(T^*L)|_{o_L}=TL\oplus T^*L$。先证明上述映射良定义。设 $X,X'$ 是沿
$o_L$ 的两个向量场，且 $[X]=[X']$ 于 $\Gamma(\nu(o_L))$ 中成立。为简便起见，下文只写 $o_L=L$。于是对所有 $x\in L$，
$X(x)-X'(x)\in T_xL$。因为 $L$ 是拉格朗日的，
\[
 (X-X')(x)\mathbin{\lrcorner}\omega_0|_{T_xL}=0,
 \quad\text{即}\quad
 X(x)\mathbin{\lrcorner}\omega_0|_{T_xL}
 =X'(x)\mathbin{\lrcorner}\omega_0|_{T_xL}.
\]
所以 $[X]\mathbin{\lrcorner}\omega_0=[X']\mathbin{\lrcorner}\omega_0$，映射良定义。类似地，它是双射；证明留给读者。

前面已经说明，与 $1$-形式 $\alpha\in\Gamma(T^*L)$ 对应的
$\operatorname{Sub}(T^*L)$ 的切向量，在 $o_L$ 处切于
$\operatorname{Lag}(T^*L)$，当且仅当 $\alpha$ 闭。这就建立了同构\eqref{eq:tangent-lag-closed-forms}。

对 $L$ 上给定的 $1$-形式 $\alpha$，以 $X_\alpha$ 表示沿 $o_L$ 的相伴向量场，它由
$X_\alpha\mathbin{\lrcorner}\omega_0|_{TL}=\alpha$ 定义。则变分
$d\Phi_{12}(X_\alpha)$ 在同构\eqref{eq:tangent-lag-closed-forms}下，对应于
$\Gamma(T^*L)$ 中由
$d\Phi_{12}(X_\alpha)\mathbin{\lrcorner}\omega_0|_{TL}$ 定义的元素。

为求这个截面，对每个给定向量场 $Y\in TL$，直接计算
\[
\begin{aligned}
 d\Phi_{12}(X_\alpha)\mathbin{\lrcorner}\omega_0(Y)
 &=\omega_0(d\Phi_{12}(X_\alpha),Y)\\
 &=\omega_0(X_\alpha,(d\Phi_{12})^{-1}Y)\\
 &=\omega_0(X_\alpha,Y)=\alpha(Y).
\end{aligned}
\]
这就完成证明。这里第二个等式使用 $\Phi_{12}$ 是辛的，第三个等式使用
$\Phi_{12}|_{o_L}=\operatorname{id}|_{o_L}$，最后一个等式使用 $X_\alpha$ 的定义。这就完成了证明。
\end{proof}

用这个引理证明切空间 $T_L\operatorname{Lag}(M,\omega)$ 的如下典范表示。

\begin{proposition}\label{prop:tangent-lag-closed-forms}
设 $\operatorname{Lag}_L(M,\omega)$ 是与 $L$ 具有相同拓扑类型的拉格朗日子流形集合，则有典范同构
\[
 T_L\operatorname{Lag}(M,\omega)\cong Z^1(L).
\]
\end{proposition}

\begin{proof}
因为 $L$ 是拉格朗日的，辛形式 $\omega$ 给出自然同构
\[
 \widetilde\nu(L)\longrightarrow(TL)^*=T^*L;\qquad
 [v]\longmapsto\widetilde\omega(v)|_{TL}.
\]
因此有自然同构
$T_L\operatorname{Sub}(M)\cong\Gamma(T^*L)$。另一方面，给定任意余切丛坐标卡，$L$ 附近的
$\operatorname{Lag}_L(M,\omega)$ 由 $L$ 上的 $C^1$-小闭 $1$-形式集合参数化，故其切空间也由 $Z^1(L)$ 参数化。\Cref{lem:normal-cotangent-identification}说明，切空间的参数化与余切丛坐标卡的选取无关。这就完成了证明。
\end{proof}

切丛
$T\operatorname{Lag}(M,\omega)\to\operatorname{Lag}(M,\omega)$ 含有一个典范分布：在
$L\in\operatorname{Lag}(M,\omega)$ 处，它由正合 $1$-形式集合
$B^1(L)\subset Z^1(L)\cong T_L\operatorname{Lag}(M,\omega)$ 张成。以
$\mathcal E\to\operatorname{Lag}(M,\omega)$ 表示相应子丛。它定义了余秩
$b^1(L)=\dim H^1(L,\mathbb R)$ 的分布。

下面这个在 (Wn90) 中陈述的定理，将在\cref{sec:exact-isotopy-hamiltonian-isotopy}中证明。

\begin{theorem}\label{thm:hamiltonian-leaves}
分布 $\mathcal E\subset T\operatorname{Lag}(M,\omega)$ 是可积子丛；穿过给定拉格朗日子流形
$L\in\operatorname{Lag}(M,\omega)$ 的叶
$\mathcal L_L\subset\operatorname{Lag}(M,\omega)$ 由 $L$ 的哈密顿形变组成。更确切地说，对每个 $L'\in\mathcal L_L$，存在
$\phi\in\operatorname{Ham}(M,\omega)$ 使 $L'=\phi(L)$。事实上，
$\operatorname{Ham}(M,\omega)$ 作用在 $\operatorname{Lag}(M,\omega)$ 上，而
$\mathcal E_L$ 是 $\operatorname{Ham}(M,\omega)$ 在 $L$ 上轨道的切空间。
\end{theorem}
